\section{Same-host capacity-configuration validation}
\label{sec:capacity-validation}

\subsection{Frozen design and accounting}
Two complementary studies test whether replay labels and candidate choices remain reliable
when allocated capacity and input multiplicity change together. The first pairs
$N=2$ and $N=4$ on one host with four A800 80GB PCIe GPUs, using
Qwen2.5-7B-Instruct FP16, vLLM 0.9.0, PyTorch 2.7.0, and independent TP1
instances. It uses the previously studied recovery and middle-burst inputs,
not new held-out windows. Two-replica inputs contain 478/509 requests
(358/389 online and 120 offline). Four-replica inputs duplicate every row
with independent IDs, producing 956/1,018 requests with unchanged arrival
times, token budgets, and deadlines. This is capacity-matched weak scaling,
not a same-load speedup comparison.

The matrix is $2$ capacities $\times$ $2$ windows $\times$
$3$ policies (all$N$, H, H+guard) $\times$ $2$ physical repeats: 24 cells,
executed serially. Repeat 1 is primary; repeat 2 tests robustness.
All$N$ remains ready throughout; H uses scale-out thresholds 8/16 at
$N=2/4$, step size one, and a 60~s cooldown. The guard evaluates legal
targets $0,\ldots,N$. Local request slots and synthetic-API admission slots
are four per allocated GPU. The observation interval is 1,800~s of arrivals
plus 300~s of continued control and dispatch. Setup precedes the observation
origin; cleanup follows the cutoff. PPO is not part of this study: the
original three-logit learned policies retain their two-device action space.

Service and lifecycle calibration is separate for each capacity and precedes
formal results. Each GPU measures 71 service points, three times each;
ten loaded launch rounds provide 20/40 startup samples at $N=2/4$.
Local startup medians are 32.590/39.729~s; guard startup safety p95 values
are 79.571/42.066~s. Samples retain their observed cache conditions; no
cold-cache claim is made. Each capacity's guard profile is common to all
predictors and physical execution. Source, input, profiles, 48 complete
predictions, and selection sets are hash-bound before the first formal
setup. No scientific failure is rerun or used to retune a predictor.

All 24 cells are technically valid, with all 2,100 decision ticks present
per cell. There are two technical smoke executions and no formal retries;
these and calibration are outside the 24 scientific cells. Every matched
all$N$ capacity reference passes strict P1+, so no dynamic cell is skipped.
Strict feasibility includes every planned online and offline request,
including undispatched and unfinished work. Table~\ref{tab:capacity-physical}
reports all physical results. Operating costs exclude setup, cleanup and
miss penalties. Savings are defined only for feasible policy/reference
pairs at the same $N$, window and repeat; costs across capacities are not
subtracted to claim savings.

\begin{table*}[t]
\centering\small
\caption{All 24 capacity-study cells, grouped by their two repeats.
Offline counts are on-time completions out of 120 ($N=2$) or 240 ($N=4$)
per repeat. All online requests meet deadlines in every cell.
Costs are operating scenario USD; $\dagger$ marks an infeasible execution.
P1+ reports passes/runs; other slash-separated values are repeat 1 / repeat 2.}
\label{tab:capacity-physical}
\setlength{\tabcolsep}{3pt}
\begin{tabular*}{\textwidth}{@{\extracolsep{\fill}}rllclll@{}}
\toprule
$N$ & Window & Policy & P1+ & Offline on time & Cost & Saving vs all$N$ (\%)\\
\midrule
2 & Recovery & all$N$ & 2/2 & 120/120 & 4.2000/4.2000 & 0.00/0.00 \\
2 & Recovery & H & 2/2 & 120/120 & 1.1468/1.1452 & 72.70/72.73 \\
2 & Recovery & H+guard & 2/2 & 120/120 & 1.0733/0.9645 & 74.45/77.04 \\
\midrule
2 & Burst & all$N$ & 2/2 & 120/120 & 4.2000/4.2000 & 0.00/0.00 \\
2 & Burst & H & 2/2 & 120/120 & 0.9659/0.9658 & 77.00/77.00 \\
2 & Burst & H+guard & 2/2 & 120/120 & 1.1507/0.9660 & 72.60/77.00 \\
\midrule
4 & Recovery & all$N$ & 2/2 & 240/240 & 8.4000/8.4000 & 0.00/0.00 \\
4 & Recovery & H & 0/2 & 238/224 & 1.8817$\dagger$/1.4377$\dagger$ & --/-- \\
4 & Recovery & H+guard & 2/2 & 240/240 & 1.8444/1.5069 & 78.04/82.06 \\
\midrule
4 & Burst & all$N$ & 2/2 & 240/240 & 8.4000/8.4000 & 0.00/0.00 \\
4 & Burst & H & 1/2 & 236/240 & 1.6000$\dagger$/1.4830 & --/82.35 \\
4 & Burst & H+guard & 2/2 & 240/240 & 1.4768/1.4828 & 82.42/82.35 \\
\bottomrule
\end{tabular*}
\end{table*}

\subsection{Physical outcomes and observed mechanism}
All 12 two-replica cells pass. At four replicas, all four all$N$ references
and all four H+guard cells pass, while H passes one of four executions.
Its recovery repeats leave two and sixteen offline jobs undispatched;
its first burst repeat leaves four. Every unfinished job is already due
before the 2,100~s cutoff. There are no late completions in this study.

The first four-replica recovery H run illustrates a tail-queue mechanism:
its last active worker is shut down at 1,774.001~s, before two final offline
arrivals at 1,785~s with deadlines at 2,085~s. The queue stays below H's
scale-out threshold of 16, and capacity remains zero through drain.
The guarded recovery run instead raises its target at 2,019~s with four
offline jobs pending and subsequently meets every deadline. These are
separate physical trajectories, not a controlled single-event intervention.
Their frozen Local target predictions first differ from physical targets
at tick 330, so a later restart alone cannot explain all preceding divergence.

The hardware boundary matters. Four-replica all$N$ recovery/burst references
serve 150/165 requests on each of GPUs 2 and 3 per repeat. In all eight
four-replica dynamic cells, GPUs 2 and 3 serve zero requests. Ready capacity
of at least three GPUs is observed for about 64--67~s, largely during initial
capacity release; the guarded recovery primary also raises its target to four
near the cutoff. Target/ready counts therefore do not establish service on
the extra devices. These original windows establish capacity-configuration sensitivity and
candidate-screening consequences. The pressure study below additionally
exercises service on both extra devices after their initial release.

\subsection{Frozen prediction and selection consequences}
For each capacity and predictor, four distinct dynamic policy/window
predictions are compared with four physical labels in each repeat.
The all$N$ references are excluded from Table~\ref{tab:capacity-prediction};
all predictors accept those feasible references. Repeat comparisons reuse
the same predictions and do not form independent prediction samples.

\begin{table}[t]
\centering\small
\caption{Dynamic-pool prediction labels and frozen two-window selection.
Agreement denominators are four physical policy/window labels per repeat;
FA/FR are descriptive false-accept/reject counts, not population error rates.
Paired entries are repeat 1 and repeat 2.
Selection retains all ties.}
\label{tab:capacity-prediction}
\setlength{\tabcolsep}{3pt}
\begin{tabular*}{\linewidth}{@{\extracolsep{\fill}}rlllll@{}}
\toprule
$N$ & Replay & Agreement & FA & FR & Selected set\\
\midrule
2 & Share & 2/4; 2/4 & 0/0 & 2/2 & H+guard \\
2 & Lookup-69 & 4/4; 4/4 & 0/0 & 0/0 & \{H, H+guard\} \\
2 & Lookup-30 & 4/4; 4/4 & 0/0 & 0/0 & \{H, H+guard\} \\
2 & Local & 4/4; 4/4 & 0/0 & 0/0 & \{H, H+guard\} \\
\midrule
4 & Share & 3/4; 2/4 & 1/1 & 0/1 & H+guard \\
4 & Lookup-69 & 2/4; 3/4 & 2/1 & 0/0 & \{H, H+guard\} \\
4 & Lookup-30 & 3/4; 4/4 & 1/0 & 0/0 & H+guard \\
4 & Local & 3/4; 4/4 & 1/0 & 0/0 & H+guard \\
\bottomrule
\end{tabular*}
\end{table}

Selection requires a candidate to be predicted strictly feasible on both
windows, then minimizes equal-window mean operating cost over
$\{\mathrm{H},\mathrm{H+guard}\}$, retaining all exact cost ties.
At $N=2$, Lookup-69, Lookup-30 and Local tie both candidates, and both are
physically feasible. At $N=4$, Lookup-69 still assigns the same predicted
trajectory and cost to both policies in each window. Its selected set thus
contains H, which is infeasible across the two windows in both physical
repeats, and H+guard, which is feasible. No tie-break is invented to convert
this set into a definite failing deployment.

Share, Lookup-30 and Local instead select H+guard at $N=4$; it is the only
candidate physically feasible on both windows in each repeat, giving zero
feasible cost regret within this two-policy pool. Nevertheless, Lookup-30
and Local falsely accept the first burst H execution, even though their
recovery rejection removes H from two-window selection. In the second
burst repeat H passes. Correct screening here therefore does not imply
correct per-window labels or a universal ranking of predictors.

\subsection{CPU audit and evidence scope}
The accompanying \path{artifact/capacity_transfer/verify.py} uses only the
Python standard library. It recomputes all 24 strict labels, clipped
occupied-resource totals, operating costs, per-GPU request counts and
48 prediction labels/costs, then reconstructs the frozen selected sets.
The records contain request outcomes, lifecycle intervals, control ticks,
and predicted state transitions; prompts and machine identifiers are removed.
This supports reduced-record audit, not model-inference reruns. Full native
receipts and frozen-source provenance are retained separately. These 24
executions and the following 12 pressure executions are additional to the
historical 197-run corpus and the eight full-day executions.

\subsection{Capacity pressure with service on all four replicas}
\label{sec:capacity-pressure}
The pressure study tests the same contract when dynamic control activates
additional replicas to serve requests. It retains every online arrival from the
original recovery input and replicates each offline job twelve times only
in $[300,1500)$~s; other offline arrivals are unchanged. Four-replica inputs
duplicate each two-replica row. Only request and window IDs change: arrival
times, token IDs, output budgets, provenance, and deadlines remain fixed.
Two-replica inputs contain 358 online and 1,000 offline jobs (1,358 total);
four-replica inputs contain 716 and 2,000 (2,716 total).

The replication factor is chosen before any new GPU observation using the
previously frozen per-capacity profiles. A recorded CPU-only protocol tests
factors $8,12,16,20,24,28,32$ in ascending order and selects the first for
which Local predicts feasible all$N$ at both capacities and N4 H+guard
predicts the required extra-card exposure. In $[300,1800)$~s, each extra
GPU must serve at least 20 requests, including at least ten after its second
or later launch, and at least three replicas must be ready for 120~s.
Factor 8 fails fourth-card exposure; factor 12 meets the criteria.
An earlier whole-input replication design fails CPU all$N$ capacity checks
and receives no GPU execution. Neither design uses new physical outcomes,
false-accept/reject counts, predicted winners, or savings to select its load.
The resulting workload is a constructed capacity-pressure test, not a
held-out sample of replay accuracy or a new production trace.

The 12-cell matrix is two capacities $\times$ three policies $\times$ two
repeats, with one fixed pressure window. It uses the same host, ordered GPU
identities, model files, software, runtime, policies, and raw-verified per-$N$
calibration as the first study. There is no new calibration or policy tuning.
Physical runs are serial, with capacities interleaved in a fixed order and
each matched all$N$ reference preceding its dynamic candidates. The
1,800~s arrival interval, 300~s drain, strict deadline contract, and
accounting rules are unchanged. All 24 complete replay predictions and
single-window selected sets are frozen before the first formal setup.
All 12 cells are technically valid, with all 2,100 control ticks present;
two technical smokes pass and no retry is used.

\begin{table*}[t]
\centering\small
\caption{All 12 pressure-study cells. Offline on-time denominators are
1,000 at $N=2$ and 2,000 at $N=4$ per repeat; every online request meets
its deadline. P1+ reports passes/runs; other slash-separated entries are
repeat 1 / repeat 2.
Costs are operating scenario USD; $\dagger$ denotes infeasibility.}
\label{tab:pressure-physical}
\setlength{\tabcolsep}{3pt}
\begin{tabular*}{\textwidth}{@{\extracolsep{\fill}}rlclll@{}}
\toprule
$N$ & Policy & P1+ & Offline on time & Cost & Saving vs all$N$ (\%)\\
\midrule
2 & all$N$ & 2/2 & 1000/1000 & 4.2000/4.2000 & 0.00/0.00 \\
2 & H & 1/2 & 996/1000 & 1.8911$\dagger$/2.0897 & --/50.24 \\
2 & H+guard & 2/2 & 1000/1000 & 1.9824/1.9856 & 52.80/52.72 \\
\midrule
4 & all$N$ & 2/2 & 2000/2000 & 8.4000/8.4000 & 0.00/0.00 \\
4 & H & 2/2 & 2000/2000 & 3.6033/3.4588 & 57.10/58.82 \\
4 & H+guard & 2/2 & 2000/2000 & 3.4779/3.7848 & 58.60/54.94 \\
\bottomrule
\end{tabular*}
\end{table*}

Table~\ref{tab:pressure-physical} reports all results. The four all$N$
references and all four H+guard executions pass. H leaves four offline
jobs unfinished in the primary two-replica run and passes its other three
executions. The second two-replica H run is feasible, so its first failure
is not treated as deterministic. No completed request is late.

All four N4 dynamic executions release and later restart both extra GPUs,
which serve 302--332 and 54--68 requests in $[300,1800)$~s after their
second or later launch (Table~\ref{tab:pressure-exposure}). Each passes the
exposure criteria. This is request-serving capacity beyond setup or a
target-only scale-out. Both policies satisfy every deadline at N4.
H+guard's feasible savings against matched all4 are 58.60\% and 54.94\%;
these are conditional on the declared scenario ledger, not cross-capacity
speedup or real billing savings.

\begin{table}[t]
\centering\small
\caption{Observed N4 capacity exposure in $[300,1800)$~s. GPU counts
include only requests dispatched after that device's second or later launch.
Ready means active and accepting requests, excluding draining.}
\label{tab:pressure-exposure}
\setlength{\tabcolsep}{4pt}
\begin{tabular*}{\linewidth}{@{\extracolsep{\fill}}lrrrr@{}}
\toprule
Policy & Repeat & GPU 2 & GPU 3 & Ready $\geq3$ (s)\\
\midrule
H & 1 & 332 & 62 & 328.61 \\
H & 2 & 330 & 54 & 340.49 \\
H+guard & 1 & 302 & 65 & 336.52 \\
H+guard & 2 & 303 & 68 & 335.59 \\
\bottomrule
\end{tabular*}
\end{table}

\subsection{Pressure-study selection and cost robustness}
Selection in this study requires predicted P1+ on its single window, then
minimizes predicted operating cost among H and H+guard, retaining exact
ties; an empty eligible set yields abstention. It is separate from the
original study's two-window selection. Each capacity/predictor has two
distinct dynamic predictions, reused for the two physical repeats;
all$N$ remains a capacity reference outside selection.

\begin{table}[t]
\centering\small
\caption{Pressure-study dynamic predictions and frozen single-window selection.
Agreement uses two physical labels per repeat; FA/FR are counts, not
estimated population rates. Paired entries are repeat 1 and repeat 2.}
\label{tab:pressure-prediction}
\setlength{\tabcolsep}{3pt}
\begin{tabular*}{\linewidth}{@{\extracolsep{\fill}}rlllll@{}}
\toprule
$N$ & Replay & Agreement & FA & FR & Selected set\\
\midrule
2 & Share & 2/2; 1/2 & 0/0 & 0/1 & H+guard \\
2 & Lookup-69 & 1/2; 2/2 & 1/0 & 0/0 & \{H, H+guard\} \\
2 & Lookup-30 & 2/2; 1/2 & 0/0 & 0/1 & H+guard \\
2 & Local & 2/2; 1/2 & 0/0 & 0/1 & H+guard \\
\midrule
4 & Share & 2/2; 2/2 & 0/0 & 0/0 & H+guard \\
4 & Lookup-69 & 0/2; 0/2 & 0/0 & 2/2 & abstain \\
4 & Lookup-30 & 1/2; 1/2 & 0/0 & 1/1 & H+guard \\
4 & Local & 2/2; 2/2 & 0/0 & 0/0 & H \\
\bottomrule
\end{tabular*}
\end{table}

At N4, Lookup-69 rejects both candidates and abstains, although both pass
both physical repeats. In contrast, it accepts and ties both on the original
short windows, where H is infeasible in three of four executions.
The same startup proxy can therefore produce false acceptance in one
capacity/load condition and false rejection in another.
Share and Local match all four N4 pressure labels; Lookup-30 falsely
rejects both H executions and selects the feasible guard.
This agreement does not establish a universal best predictor.

Local selects H at N4, while Share and Lookup-30 select H+guard.
In the primary repeat, H costs 3.6033 versus H+guard's 3.4779 scenario USD,
so Local's feasible cost gap is 3.61\%. In the second repeat, H costs
3.4588 versus 3.7848 and is the cheaper candidate; selecting H+guard then
incurs a 9.43\% gap. These within-repeat gaps describe the cost consequence
of each frozen choice. The physical ranking reverses between repeats;
with two observations per candidate, this does not isolate a systematic replay
ranking error or establish either policy as consistently cheaper.
At N2, Local, Lookup-30, and Share select H+guard, which is the cheaper
feasible candidate in both repeats; Lookup-69's tie includes the failed
primary H execution.

The prompt-free \path{artifact/capacity_pressure/verify.py} recomputes all
12 physical labels and costs, per-GPU restart exposure, and 24 frozen
prediction labels/costs and selected sets using the standard library.
It preserves all planned requests and both repeats. Together, the two
studies add 36 physical cells (32 feasible and four infeasible) and 72
frozen predictions, with separate input-selection and screening rules.
The original-window and constructed-pressure results are reported
separately rather than pooled into a population accuracy rate.
